\documentclass[12pt,a4paper]{article}

\usepackage[utf8]{inputenc}
\usepackage[T2A]{fontenc}
\usepackage[russian]{babel}
\usepackage{amsmath,amssymb}
\usepackage{geometry}
\usepackage{graphicx}
\usepackage{float}

\geometry{margin=2.2cm}

\newcommand{\plotplaceholder}[2][0.9\textwidth]{%
    \IfFileExists{#2}{%
        \includegraphics[width=#1]{#2}%
    }{%
        \fbox{%
            \begin{minipage}[c][0.22\textheight][c]{#1}
            \centering
            Заглушка для графика\\[4pt]
            \texttt{\detokenize{#2}}
            \end{minipage}%
        }%
    }%
}

\begin{document}

\title{Гистограммная оценка}
\author{}
\date{}
\maketitle

\section*{Постановка}

Пусть, как и раньше,
\[
a=\sqrt{3},\qquad
f(x)=\frac{1}{2a}\mathbf{1}_{\{|x|\le a\}}.
\]
Обычная гистограмма --- это способ изобразить выборку: ось разбивается на
интервалы, и для каждого интервала показывается, сколько наблюдений туда
попало. Такая картинка зависит от конкретной выборки.

Гистограммная оценка плотности использует ту же идею, но нормирует высоты
столбцов так, чтобы получилась именно оценка плотности вероятности. Если в
разряд ширины $h$ попало $n_j$ наблюдений из $n$, то высота оценки на этом
разряде равна $n_j/(nh)$. Поэтому площадь столбца равна
\[
h\cdot\frac{n_j}{nh}=\frac{n_j}{n},
\]
то есть относительной частоте попадания в этот разряд. Суммарная площадь
всех столбцов равна $1$, как и должно быть у плотности.

В этом разделе мы изучаем не одну случайную картинку-гистограмму, а среднее
качество гистограммной оценки как метода. Для этого считаем математическое
ожидание интегральной квадратической ошибки и затем нормируем его так же,
как в разделе про ядерную оценку.

Рассмотрим упрощённый вариант гистограммной оценки. Фиксируем
рассматриваемый отрезок $[-a,a]$ и выбираем числовой параметр $h>0$ ---
ширину разряда гистограммы. Границы разрядов задаются формулой
\[
x_j=-a+jh.
\]
Если последняя точка такого равномерного шага выходит за правую границу
носителя, то последний разряд обрезается в точке $a$. Поэтому фактические
разряды имеют вид
\[
I_j=(x_j,x_{j+1}),\qquad j=0,\ldots,m-1,
\]
где $x_0=-a$, $x_m=a$, а для промежуточных точек $x_j=-a+jh$.
Обозначим ширину $j$-го разряда через
\[
w_j=x_{j+1}-x_j.
\]
Для всех внутренних разрядов $w_j=h$, а последний разряд может иметь
меньшую ширину. В этой постановке $h$ не является сдвигом сетки и не
связан с параметром сглаживания ядерной оценки: это самостоятельная ширина
разряда гистограммы.

Пусть $n_j$ --- число наблюдений, попавших в интервал $I_j$.
Гистограммная оценка имеет вид
\[
\widehat f^H_h(x)=\frac{n_j}{nw_j},
\qquad x\in I_j.
\]
В точках разрыва можно задать любое конечное значение: на интегральную
квадратическую ошибку это не влияет.

\section*{Точная формула для MISE}

Чтобы посчитать среднюю ошибку гистограммы, нужно знать не сами случайные
числа $n_j$, а вероятности попадания наблюдения в разряды. Обозначим такую
вероятность через $p_j$:
\[
p_j=\mathbb P\{X\in I_j\}
=\int_{I_j} f(x)\,dx.
\]
Тогда $n_j$ --- случайная величина с биномиальным распределением: в каждом
из $n$ независимых наблюдений есть вероятность $p_j$ попасть в данный
разряд.

Для равномерной плотности на выбранном отрезке это просто нормированная
ширина разряда:
\[
p_j=\frac{w_j}{2a},
\qquad
w_j=\operatorname{length}(I_j).
\]
Сумма по $j$ конечна, и $\sum_j p_j=1$.

Теперь раскладываем ошибку
\[
\operatorname{MISE}_H
=
\mathbb E\int\bigl(\widehat f^H(x)-f(x)\bigr)^2\,dx
\]
по формуле квадрата:
\[
\mathbb E\int\bigl(\widehat f^H-f\bigr)^2
=
\int f^2
-2\int \mathbb E\widehat f^H\, f
+\int \mathbb E(\widehat f^H)^2.
\]
Первые два слагаемых отвечают за смещение оценки, а последнее содержит
дисперсионный вклад из-за случайности выборки.

Найдём смешанное слагаемое. Внутри одного разряда математическое ожидание
оценки равно $p_j/w_j$, потому что $\mathbb E n_j=np_j$. Поэтому
\[
\int \mathbb E\widehat f^H(x)f(x)\,dx
=\sum_j \frac{p_j}{w_j}\int_{I_j} f(x)\,dx
=\sum_j \frac{p_j^2}{w_j}.
\]

Теперь найдём второе среднее слагаемое. Так как $n_j$ имеет биномиальное
распределение с параметрами $n$ и $p_j$,
\[
\mathbb E n_j^2=np_j(1-p_j)+n^2p_j^2.
\]
Следовательно,
\[
\begin{aligned}
\int \mathbb E\bigl(\widehat f^H(x)\bigr)^2\,dx
&=\sum_j w_j\,\frac{\mathbb E n_j^2}{n^2w_j^2}\\
&=\sum_j\frac{p_j(1-p_j)}{nw_j}
+\sum_j\frac{p_j^2}{w_j}.
\end{aligned}
\]
Подставляя найденные выражения в разложение квадрата ошибки, получаем
точную формулу:
\[
\begin{aligned}
\operatorname{MISE}_H
&=\mathbb E\int\bigl(\widehat f^H(x)-f(x)\bigr)^2\,dx\\
&=\int f^2(x)\,dx
-2\int \mathbb E\widehat f^H(x)f(x)\,dx
+\int \mathbb E\bigl(\widehat f^H(x)\bigr)^2\,dx\\
&=\|f\|_2^2
-\sum_j\frac{p_j^2}{w_j}
+\sum_j\frac{p_j(1-p_j)}{nw_j}.
\end{aligned}
\]

\section*{Нормированная ошибка}

Так как
\[
\|f\|_2^2=\int_{-a}^{a}\frac{1}{4a^2}\,dx=\frac{1}{2a},
\]
нормированная ошибка равна
\[
\operatorname{RMISE}_H(h)
=
\frac{\operatorname{MISE}_H}{\|f\|_2^2}
=
2a\left(
\|f\|_2^2
-\sum_j\frac{p_j^2}{w_j}
+\sum_j\frac{p_j(1-p_j)}{nw_j}
\right).
\]
В нашем равномерном примере $p_j=w_j/(2a)$, поэтому первые два слагаемых
сокращаются:
\[
\|f\|_2^2-\sum_j\frac{p_j^2}{w_j}
=
\frac{1}{2a}-\sum_j\frac{w_j}{4a^2}
=0.
\]
Остаётся только дисперсионная часть
\[
\operatorname{RMISE}_H(h)
=
2a\sum_j\frac{p_j(1-p_j)}{nw_j}.
\]
Если разбиение состоит из $m$ разрядов, то для такого прямоугольного закона
получается простая формула
\[
\operatorname{RMISE}_H(h)=\frac{m-1}{n},
\qquad
m=\left\lceil\frac{2a}{h}\right\rceil .
\]
Это специальное свойство выбранного закона: математическое ожидание
гистограммы совпадает с истинной плотностью на каждом разряде, а ошибка
возникает только из-за случайных частот в разрядах.

\section*{Примеры самих гистограмм}

Построим обычные гистограммы по одной фиксированной выборке из равномерного
распределения. Эти графики выглядят как привычные гистограммы: столбцы
показывают нормированные частоты попадания в разряды, а красная линия
показывает истинную плотность $f(x)$.

В первом наборе графиков параметром является именно ширина разряда $h$.
Число разрядов при этом не задаётся отдельно, а получается из выбранного
шага:
\[
m=\left\lceil\frac{2a}{h}\right\rceil .
\]

\begin{figure}[H]
\centering
\plotplaceholder[0.92\textwidth]{sample_histograms_h_grid.png}
\caption{Гистограммные оценки по одной выборке при разных ширинах разряда $h$.}
\end{figure}

Для сопоставления можно показать ту же идею через число разрядов $m$.
Это не отдельная постановка: при фиксированном отрезке $[-a,a]$ параметр
$m$ однозначно задаёт шаг $h=2a/m$.

\begin{figure}[H]
\centering
\plotplaceholder[0.92\textwidth]{sample_histograms_m_grid.png}
\caption{Гистограммные оценки по одной выборке при разных числах разрядов $m$.}
\end{figure}

\section*{Графики ошибки}

График ниже уже не является самой гистограммой. Это график качества
гистограммной оценки: значение $\operatorname{RMISE}_H(h)$ считается по
точной конечной сумме вероятностей $p_j$, без моделирования выборок.

\begin{figure}[H]
\centering
\plotplaceholder[0.82\textwidth]{hist_rmise_vs_h.png}
\caption{Зависимость нормированной ОИСКО гистограммной оценки от ширины разряда $h$.}
\end{figure}

Чем меньше $h$, тем больше разрядов попадает в отрезок $[-a,a]$, поэтому
растёт дисперсионная часть ошибки. При $h\ge 2a$ весь носитель попадает в
один разряд, и в данном специальном равномерном примере оценка совпадает с
истинной плотностью почти всюду.

\end{document}
